\chapter{Primitive Roots and the Multiplicative Group}

\begin{center}\itshape
The decimal expansion of $\frac17=0.\overline{142857}$ repeats every $6$ digits, and $\frac1{17}$ every $16$.\\
Why do some primes $p$ give the longest possible period $p-1$, and are there infinitely many of them? Nobody knows.
\end{center}

In 1927 Emil Artin, chatting with Helmut Hasse, made a guess that sounds like schoolwork: the decimal expansion of $\frac1p$ has the longest possible
period, $p-1$ digits, for infinitely many primes $p$, and in fact for about $37\%$ of them. Nearly a century later nobody has proved it for the number
$10$, or for any other single base. The best result, by Heath-Brown in 1986, says that at least one of $2$, $3$, $5$ works, without saying which. The
question is about \emph{primitive roots}: elements whose powers run through every non-zero residue modulo $p$. This chapter proves that they always
exist modulo a prime, and finds out exactly which moduli have them.

\begin{reminder}[A change of coordinate]
The additive group of remainders is generated by $1$. This chapter asks when
multiplication admits an equally simple coordinate: a unit $g$ whose powers
visit all the other units. The answer depends on the modulus. For a prime it is
always yes; modulo $8$ it is already no. We will separate existence, finding a
generator, and computing the exponent of a given unit. They are three different
problems, with three very different costs.
\end{reminder}

\section{Orders}

\begin{definition}
Let $\gcd(a,n)=1$. The \emph{order} $\operatorname{ord}_n(a)$ is the least $k\ge1$ with $a^k\equiv1\pmod n$.
\end{definition}

\begin{proposition}\label{prop:order}
Let $k=\operatorname{ord}_n(a)$. Then $a^m\equiv1\pmod n$ if and only if $k\mid m$. In particular $k\mid\varphi(n)$. Moreover $\operatorname{ord}_n(a^j)=\dfrac k{\gcd(k,j)}$.
\end{proposition}
\begin{proof}
$\{m\in\Z:a^m\equiv1\}$ is a subgroup of $\Z$ containing $k$ as its least positive element, so it is $k\Z$ (Chapter~1). Euler gives $a^{\varphi(n)}\equiv1$. Finally $(a^j)^m\equiv1$ iff $k\mid jm$ iff $\frac k{\gcd(k,j)}\mid m$.
\end{proof}

The definition makes sense because Euler's theorem supplies at least one
positive exponent returning to $1$. Negative powers in the proof are powers
of the multiplicative inverse, which exists precisely because $\gcd(a,n)=1$.
For $j=0$, the order formula gives $k/\gcd(k,0)=1$, as it should: $a^0=1$.

\paragraph{The first return controls every return.}
Modulo $13$, the powers of $3$ are $1,3,9,1,3,9,\dots$, so its order is $3$.
If $3^{100}\equiv3\pmod{13}$, this is not a coincidence:
$100=3\cdot33+1$. In general, division $m=qk+r$, $0\le r<k$, gives
$a^m\equiv a^r$; the least-return property forces $r=0$ exactly when the
power is $1$. This is an alternative to the subgroup argument above.

Here is a useful distinction. Euler says $a^{\varphi(n)}\equiv1$; he does
not say that $\varphi(n)$ is the \emph{first} such exponent. Modulo $8$,
$3^2\equiv1$ although $\varphi(8)=4$. An order is a period, not a deadline.

A \emph{primitive root} modulo $n$ is an $a$ with $\operatorname{ord}_n(a)=\varphi(n)$, that is, a generator of the group $(\Z/n\Z)^\times$, which is then cyclic.

\section{\texorpdfstring{$(\Z/p\Z)^\times$}{(Z/pZ)*} is cyclic}

\begin{think}
List the powers of $2$ modulo $13$. Do they reach every non-zero class? Then compute the order of each of $1,\dots,12$ modulo $13$ and count how many elements have order $1,2,3,4,6,12$. Compare these counts with $\varphi(d)$. Can you guess the general statement, and see why a polynomial of degree $d$ having at most $d$ roots might be the key?
\end{think}

\begin{theorem}[Gauss]\label{thm:cyclic}
For $p$ prime, $(\Z/p\Z)^\times$ is cyclic: there are exactly $\varphi(p-1)$ primitive roots modulo $p$.
\end{theorem}
\begin{proof}
For $d\mid p-1$, let $N(d)$ be the number of elements of order $d$. Suppose $N(d)>0$, and let $a$ have order $d$. Then $1,a,\dots,a^{d-1}$ are $d$ distinct roots of $X^d-1$ in the field $\Z/p\Z$, and a polynomial of degree $d$ has at most $d$ roots, so these are all the elements $x$ with $x^d=1$. Those of order exactly $d$ are the $a^j$ with $\gcd(j,d)=1$ (Proposition~\ref{prop:order}): $N(d)=\varphi(d)$.
So $N(d)\le\varphi(d)$ for every $d\mid p-1$. But every element has some order dividing $p-1$, so $\sum_{d\mid p-1}N(d)=p-1=\sum_{d\mid p-1}\varphi(d)$ (Gauss' formula, Chapter~4). Hence $N(d)=\varphi(d)$ for all $d$, in particular $N(p-1)=\varphi(p-1)\ge1$.
\end{proof}

\paragraph{Why the field hypothesis matters.}
The polynomial root bound used in the proof follows by repeated division by
$X-a$: if $a$ is a root, write $P=(X-a)Q$; any other root $b$ satisfies
$(b-a)Q(b)=0$, and since $b-a\ne0$ is invertible in a field, $Q(b)=0$.
Induction on the degree gives the bound. Over $\Z/8\Z$, the polynomial
$X^2-1$ has the four roots $1,3,5,7$: a degree-two polynomial has staged a
small rebellion. The failure is exactly that non-zero factors can have product
zero in this ring. Thus the counting proof is exploiting a field property,
not merely a clever arrangement of residues.

There is also a logical point in the argument: we did \emph{not} assume that
an element of order $d$ exists for every divisor $d$. We first proved that
$N(d)$ is either $0$ or $\varphi(d)$. The equality of the two total sums then
forces every deficit to vanish. In particular the largest order exists.
For $p=2$ the same argument includes the one-element group, with
$\varphi(1)=1$ and primitive root $1$.

\begin{figure}[ht]
\centering
\fig{ch05-orders}
\caption{Orders of units modulo $13$. For each divisor $d$ of $12$, the bar
height is $N(d)=\varphi(d)$: $1,1,2,2,2,4$. The twelve units are partitioned
by their \emph{exact} orders; merely knowing $x^d=1$ counts units whose order
divides $d$, which is a different question. The rightmost bar gives the four
primitive roots.}
\label{fig:ordercounts}
\end{figure}

Figure~\ref{fig:primroot} draws the theorem for $p=13$: once the units are placed around a circle in the order of the powers of a primitive root, multiplication becomes rotation.

\begin{figure}[ht]
\centering
\fig{ch05-primroot}
\caption{The twelve units modulo $13$, placed in the order $2^0,2^1,\dots,2^{11}$. Left: multiplying by the primitive root $2$ moves one step and visits every unit. Right: $3=2^4$, so multiplying by $3$ jumps four steps, and the units split into four cycles of length $\operatorname{ord}_{13}(3)=3$. The shaded units $2,6,11,7$ are the $\varphi(12)=4$ primitive roots, the powers $2^j$ with $\gcd(j,12)=1$.}
\label{fig:primroot}
\end{figure}

\paragraph{Finding a generator without an exhausting roll call.}
Let $g$ be a unit modulo $n$ and write $M=\varphi(n)$. Then $g$ has order
$M$ if and only if
\[
 g^{M/q}\not\equiv1\pmod n
 \quad\text{for every prime }q\mid M.
\]
Indeed, if its order $d$ is a proper divisor of $M$, the integer $M/d>1$
has a prime divisor $q$, and then $d\mid M/q$. Conversely any failed test
shows that the order divides a smaller exponent. Notice that this criterion
requires the prime divisors of $M$, not a list of all its divisors.

For $n=13$, $M=12$ has prime divisors $2$ and $3$. We find
$2^6\equiv-1$ and $2^4\equiv3$ modulo $13$, so $2$ is primitive.
By contrast $4^6\equiv1$, so $4$ fails immediately; in fact its order is $6$.
This is a rigorous certificate of success or failure, and it uses just two
modular exponentiations. The hint to Exercise~\ref{ex:findroots} is now an
algorithm we have proved.

\section{Which \texorpdfstring{$(\Z/n\Z)^\times$}{(Z/nZ)*} are cyclic?}

\begin{think}
Square each of the units $1,3,5,7$ modulo $8$. What do you get? Why does this show that $(\Z/8\Z)^\times$ has no generator? Which other moduli might suffer from the same problem?
\end{think}

\begin{lemma}\label{lem:lift}
Let $p$ be an odd prime and $g$ an integer with $g^{p-1}=1+pt$, $p\nmid t$. Then for every $k\ge1$, $g^{(p-1)p^{k-1}}=1+p^kt_k$ with $p\nmid t_k$.
\end{lemma}
\begin{proof}
Induction on $k$. If $x=1+p^kt_k$ with $p\nmid t_k$, then
$x^p=1+p^{k+1}t_k+\binom p2p^{2k}t_k^2+\dots$. Every further term is divisible by $p^{k+2}$: $\binom p2p^{2k}$ because $p\mid\binom p2$ ($p$ odd) and $2k+1\ge k+2$, and the others because $jk\ge k+2$ for $j\ge3$. So $x^p=1+p^{k+1}t_{k+1}$ with $t_{k+1}\equiv t_k\not\equiv0\pmod p$.
\end{proof}

\begin{theorem}\label{thm:cyclicn}
$(\Z/n\Z)^\times$ is cyclic if and only if $n\in\{1,2,4\}$ or $n=p^k$ or $n=2p^k$ with $p$ an odd prime and $k\ge1$.
\end{theorem}
\begin{proof}
\emph{$n=p^k$.} Let $g$ be a primitive root modulo $p$. If $g^{p-1}\equiv1\pmod{p^2}$, replace $g$ by $g+p$, still primitive modulo $p$: by the binomial theorem $(g+p)^{p-1}\equiv g^{p-1}+(p-1)pg^{p-2}\equiv1-pg^{p-2}\not\equiv1\pmod{p^2}$. So we may assume $g^{p-1}=1+pt$ with $p\nmid t$. Let $m=\operatorname{ord}_{p^k}(g)$. It divides $\varphi(p^k)=(p-1)p^{k-1}$, and is a multiple of $\operatorname{ord}_p(g)=p-1$, so $m=(p-1)p^j$ with $j\le k-1$. If $j<k-1$, then $g^{(p-1)p^{k-2}}\equiv1\pmod{p^k}$, contradicting the lemma (applied with $k-1$). So $m=\varphi(p^k)$.

\emph{$n=2p^k$.} By the Chinese remainder theorem $(\Z/2p^k\Z)^\times\cong(\Z/2\Z)^\times\times(\Z/p^k\Z)^\times$, and the first factor is trivial. The cases $n=1,2,4$ are checked by hand ($3$ generates modulo $4$); at $n=1$ the unit group is understood as the one-element group.

\emph{The others are not cyclic.} A cyclic group of even order $m$, generated by $g$, has exactly two solutions of $x^2=1$: $(g^j)^2=1$ iff $m\mid2j$ iff $j\in\{0,\frac m2\}$ modulo $m$. For $n>2$, $\varphi(n)$ is even, so it suffices to exhibit four square roots of $1$. If $8\mid n$: $\pm1$ and $2^{a-1}\pm1$ modulo $2^a$ ($a\ge3$) give four, which combine with $1$ modulo the odd part by the Chinese remainder theorem. If $n$ has two distinct odd prime factors, or $n=4p^k$, then $n=ab$ with $a,b>2$ coprime, and the pairs $(\pm1,\pm1)$ give four.
\end{proof}

\paragraph{Watching a primitive root lift.}
Take $p=3$ and $g=2$. Since $2^2=1+3$, the lemma applies with $t=1$.
The orders modulo $3,9,27$ are therefore $2,6,18$:
\[
 2^2=1+3,\qquad 2^6=64=1+9\cdot7,
 \qquad 2^{18}\equiv1+27\pmod{81}.
\]
The first non-zero power of $3$ in $g^{2\cdot3^{k-1}}-1$ increases by
exactly one each time we cube. It cannot skip a level. This is what prevents
the order from stopping short of $2\cdot3^{k-1}$ modulo $3^k$.
If the first candidate has $g^{p-1}\equiv1\pmod{p^2}$, changing it to
$g+p$ repairs precisely that first step; its class modulo $p$ does not change.

The non-cyclic cases are worth comparing directly. Modulo $15$, the four
square roots of $1$ are $1,4,11,14$, the Chinese remainder combinations of
$\pm1$ modulo $3$ and $5$. Modulo $8$ they are $1,3,5,7$.
A cyclic group of even order has room for just one non-identity element whose
square is $1$. These four roots are an obstruction one can see, without
trying every possible generator.

\section{Discrete logarithms}

\begin{think}
Find $i$ with $2^i\equiv5\pmod{13}$. How did you do it? Now imagine $p$ has $600$ digits: how long would your method take, compared with computing $2^i\bmod p$ from $i$?
\end{think}

If $g$ is a primitive root modulo $p$, every unit is $g^i$ for a unique $i\in\Z/(p-1)\Z$, its \emph{index} or \emph{discrete logarithm} $\operatorname{ind}_g$. It turns multiplication into addition: $\operatorname{ind}(ab)=\operatorname{ind}a+\operatorname{ind}b\pmod{p-1}$, so multiplicative questions become linear ones modulo $p-1$ (Exercise~\ref{ex:kth}).
Computing $g^i$ from $i$ is fast (square and multiply); recovering $i$ from $g^i$ is believed to be hard in appropriately chosen large groups. A large $p$ alone is not enough: the factorisation of $p-1$ and the chosen subgroup matter, as does the algorithm available. This asymmetry is the basis of the Diffie--Hellman key exchange.

\paragraph{Solving a power equation in the exponent coordinate.}
For $x^4\equiv3\pmod{13}$, use $g=2$, so $3=2^4$ and write $x=2^u$.
We need $4u\equiv4\pmod{12}$, or $u\equiv1\pmod3$. The four exponents
$1,4,7,10$ give the four roots
\[
 x\equiv2,3,11,10\pmod{13}.
\]
Each fourth power is $3$; these are all the roots because every non-zero
class has a unique exponent modulo $12$. The exponent lives modulo $p-1$,
while the original residue lives modulo $p$. Confusing those two clocks is
the standard trap in discrete-logarithm calculations.

\begin{think}
Would the equation $x^4\equiv2\pmod{13}$ have a solution? The exponent
congruence would be $4u\equiv1\pmod{12}$, impossible because its left side
is divisible by $4$. Equivalently, $2^3\not\equiv1\pmod{13}$.
Now try the same question for $x^3\equiv5\pmod{13}$, using $5=2^9$.
The exponent condition is $u\equiv3\pmod4$, giving $x\equiv8,11,7$.
\end{think}

\section{Exercises}

\begin{exercise}[$\star$]\label{ex:findroots}
Find a primitive root modulo $7$, $11$ and $13$, and list all primitive roots modulo $13$.
\end{exercise}

\begin{exercise}[$\star$]\label{ex:period}
Let $p\ne2,5$ be prime. Show that the decimal expansion of $\frac1p$ is periodic with least period $\operatorname{ord}_p(10)$. Check the periods of $\frac17$ and $\frac1{17}$.
\end{exercise}

\begin{exercise}[$\star\star$]\label{ex:fermatfactor}
Show that every prime factor $p$ of the Fermat number $2^{2^n}+1$ satisfies $p\equiv1\pmod{2^{n+1}}$. (This is how Euler found $641\mid2^{32}+1$.)
\end{exercise}

\begin{exercise}[$\star\star$]\label{ex:powersum}
Let $p$ be prime and $k\ge0$. Show that $\displaystyle\sum_{x=1}^{p-1}x^k\equiv\begin{cases}-1&\text{if }(p-1)\mid k,\\0&\text{otherwise}\end{cases}\pmod p$.
\end{exercise}

\begin{exercise}[$\star\star$]\label{ex:kth}
Let $p$ be prime, $k\ge1$ an integer, $p\nmid a$, and $d=\gcd(k,p-1)$. Show that $x^k\equiv a\pmod p$ has a solution if and only if $a^{(p-1)/d}\equiv1\pmod p$, and that it then has exactly $d$ solutions.
\end{exercise}

\begin{exercise}[$\star\star\star$]\label{ex:gausswilson}
\emph{Gauss' generalisation of Wilson's theorem.} Let $n>2$. Show that the product of all units modulo $n$ is $\equiv-1\pmod n$ if $(\Z/n\Z)^\times$ is cyclic, and $\equiv1$ otherwise. \emph{Hint: pair each unit with its inverse, and study the subgroup $\{x:x^2=1\}$.}
\end{exercise}

\begin{exercise}[$\star\star\star$]\label{ex:cyclotomic}
\emph{Infinitely many primes $\equiv1\pmod n$.} For a positive integer $n$, define the cyclotomic polynomials by $X^n-1=\prod_{d\mid n}\Phi_d(X)$, that is, $\Phi_n=\prod_{\gcd(k,n)=1}\left(X-e^{2ik\pi/n}\right)$.
\begin{enumerate}[label=(\alph*)]
\item Show by induction that $\Phi_n$ is monic with integer coefficients, and that $\Phi_n(0)=\pm1$ for $n\ge2$.
\item Let $p$ be a prime with $p\nmid n$ and $p\mid\Phi_n(a)$. Show that $\operatorname{ord}_p(a)=n$, hence $p\equiv1\pmod n$. \emph{Hint: otherwise $a$ would be a double root of $X^n-1$ modulo $p$.}
\item Deduce that there are infinitely many primes $p\equiv1\pmod n$.
\end{enumerate}
\end{exercise}

\begin{beauty}
Theorem~\ref{thm:cyclic} is the reason finite fields are so pleasant: their multiplicative groups are cyclic, a fact Tome~2 proves for every finite field $\mathbb F_q$ by the same counting argument. \emph{Artin's conjecture} (1927) predicts that $10$ is a primitive root modulo a proportion $\approx37.4\%$ of the primes. More generally every integer other than $-1$ and perfect squares should be primitive for a positive proportion of primes, but that proportion can depend on the integer: $37.4\%$ is not a universal answer. For a specified base such as $10$ the assertion remains open unconditionally; Hooley proved it under the generalised Riemann hypothesis for the relevant number fields (1967). Exercise~\ref{ex:cyclotomic} is the special case $a\equiv1$ of Dirichlet's theorem on primes in progressions, proved in full in Tome~5.
\end{beauty}

\section{Hints}

\begin{hint}{ex:findroots}
$g$ is a primitive root modulo $p$ if and only if $g^{(p-1)/q}\not\equiv1$ for every prime $q\mid p-1$. Once you have one primitive root $g$, the others are the $g^j$ with $\gcd(j,p-1)=1$.
\end{hint}
\begin{hint}{ex:period}
In the long division of $1$ by $p$, the remainder after $j$ steps is $10^j\bmod p$. When does it come back to $1$?
\end{hint}
\begin{hint}{ex:fermatfactor}
From $2^{2^n}\equiv-1\pmod p$, find the exact order of $2$ modulo $p$, then use that it divides $p-1$.
\end{hint}
\begin{hint}{ex:powersum}
Write the units as $g^0,\dots,g^{p-2}$ for a primitive root $g$. The sum becomes geometric with ratio $g^k$.
\end{hint}
\begin{hint}{ex:kth}
Take discrete logarithms: $x=g^\xi$, $a=g^\alpha$. The equation becomes $k\xi\equiv\alpha\pmod{p-1}$, a linear congruence from Chapter~3.
\end{hint}
\begin{hint}{ex:gausswilson}
After pairing each unit with its inverse, only the solutions of $x^2=1$ remain. How many are there when the group is cyclic, and what is their product when there are $2^r\ge4$ of them?
\end{hint}
\begin{hint}{ex:cyclotomic}
(a) Divide $X^n-1$ by $\prod_{d\mid n,\,d<n}\Phi_d$, a monic polynomial with integer coefficients. (b) If $\operatorname{ord}_p(a)=d<n$, then modulo $p$ the class of $a$ is a root of $X^d-1$, hence of some $\Phi_e$ with $e\mid d$, and also of $\Phi_n$: a double root of $X^n-1$. Look at the derivative. (c) If $p_1,\dots,p_r$ were all the primes $\equiv1\pmod n$, evaluate $\Phi_n$ at a large multiple of $m=np_1\cdots p_r$: the value is $\equiv\pm1\pmod m$.
\end{hint}

\section{Solutions}

\begin{solution}{ex:findroots}
\idea To test whether $g$ is primitive modulo $p$, we need not compute all its powers: by Proposition~\ref{prop:order} its order divides $p-1$, so it is primitive iff $g^{(p-1)/q}\not\equiv1$ for each prime $q\mid p-1$. Try $g=2,3,\dots$. Once one primitive root is known, all of them are its powers with exponent prime to $p-1$.
\writeup Modulo $7$: powers of $3$ are $3,2,6,4,5,1$, so $3$ is primitive. Modulo $11$: $2$ has order dividing $10$, and $2^2=4$, $2^5=32\equiv-1$, so the order is $10$. Modulo $13$: $2^4=16\equiv3$, $2^6=64\equiv-1$, so $\operatorname{ord}(2)$ divides $12$ but not $4$ or $6$: it is $12$. The primitive roots modulo $13$ are $2^j$ with $\gcd(j,12)=1$ (Proposition~\ref{prop:order}): $2^1,2^5,2^7,2^{11}\equiv2,6,11,7$. There are $\varphi(12)=4$.
\end{solution}

\begin{solution}{ex:period}
\idea In long division, the digits are determined by the successive remainders, and the remainder after $j$ steps is $10^j\bmod p$. So the expansion repeats exactly when the powers of $10$ modulo $p$ do, i.e.\ with period the order of $10$.
\writeup Long division of $1$ by $p$ produces the remainders $r_j=10^j\bmod p$ (each step multiplies the remainder by $10$ and reduces), and the $(j+1)$-th digit is $\lfloor10r_j/p\rfloor$, a function of $r_j$. Since $10$ is invertible modulo $p$, the sequence $10^j\bmod p$ is purely periodic with least period $\operatorname{ord}_p(10)$, and so are the digits; the least period of the digits is not smaller, because $\frac1p=\frac{A}{10^k-1}$ with $A$ an integer (period $k$) means exactly $p\mid10^k-1$. Modulo $7$, $10\equiv3$ is primitive: period $6$. Modulo $17$: $10^2\equiv-2$, $10^8\equiv16\equiv-1$, so the order is $16$.
\end{solution}

\begin{solution}{ex:fermatfactor}
\idea $p\mid2^{2^n}+1$ says $2^{2^n}\equiv-1$, hence $2^{2^{n+1}}\equiv1$: we know the order of $2$ divides a power of $2$, and $-1\ne1$ pins it down exactly. Then order divides $p-1$.
\writeup $2^{2^n}\equiv-1\pmod p$, so $2^{2^{n+1}}\equiv1$. The order of $2$ divides $2^{n+1}$ but not $2^n$ (else $2^{2^n}\equiv1\not\equiv-1$, as $p$ is odd). So $\operatorname{ord}_p(2)=2^{n+1}$, which divides $p-1$.
\end{solution}

\begin{solution}{ex:powersum}
\idea With a primitive root $g$, the sum over all units becomes a geometric sum $\sum_i(g^k)^i$, which is easy: it depends on whether the ratio $g^k$ is $1$.
\writeup Let $g$ be a primitive root. The $x$ run over $g^0,\dots,g^{p-2}$, so the sum is $\sum_{i=0}^{p-2}(g^k)^i$. If $(p-1)\mid k$, $g^k\equiv1$ and the sum is $p-1\equiv-1$. Otherwise $g^k\not\equiv1$ and the geometric sum is $\frac{g^{k(p-1)}-1}{g^k-1}\equiv0$, since $g^{k(p-1)}\equiv1$ and $g^k-1$ is invertible.
\end{solution}

\begin{solution}{ex:kth}
\idea The discrete logarithm turns the multiplicative equation $x^k=a$ into the linear congruence $k\xi\equiv\alpha\pmod{p-1}$, which Chapter~3 solves completely. Then translate the condition $d\mid\alpha$ back into a condition on $a$.
\writeup Write $a=g^\alpha$ and $x=g^\xi$. The equation becomes $k\xi\equiv\alpha\pmod{p-1}$, solvable if and only if $d\mid\alpha$, with exactly $d$ solutions modulo $p-1$ (Proposition~\ref{prop:lincong}). And $d\mid\alpha$ iff $(p-1)\mid\alpha\frac{p-1}d$ iff $a^{(p-1)/d}=g^{\alpha(p-1)/d}\equiv1$.
\end{solution}

\begin{solution}{ex:gausswilson}
\idea Copy the proof of Wilson's theorem: pairing $x$ with $x^{-1}$ reduces the product to the product of the self-inverse elements, the solutions of $x^2=1$. These form a group where every element has order $\le2$, which is a vector space over $\mathbb F_2$; the sum of all vectors of such a space is $0$ unless it is a line.
\writeup Pairing each unit with its inverse, all pairs $\{x,x^{-1}\}$ with $x\ne x^{-1}$ contribute $1$, so the product equals the product of the elements of $S=\{x:x^2=1\}$. $S$ is a subgroup in which every element is its own inverse. If $|S|=2$, $S=\{1,-1\}$ and the product is $-1$. Otherwise $|S|>2$. Written additively, $S$ is an abelian group in which every element has order at most $2$, that is, a vector space over $\mathbb F_2$, of some dimension $r\ge2$. The sum of all vectors of $\mathbb F_2^r$ is $0$ when $r\ge2$, because each coordinate equals $1$ in exactly $2^{r-1}$ vectors, an even number. So the product is $1$. Finally, by the proof of Theorem~\ref{thm:cyclicn}, $|S|=2$ exactly when $(\Z/n\Z)^\times$ is cyclic.
\end{solution}

\begin{solution}{ex:cyclotomic}
\idea This is Euclid's argument with $\Phi_n$ in place of $X+1$: we need a polynomial whose prime factors (apart from those of $n$) are all $\equiv1\pmod n$. Part (b) gives this: a prime dividing $\Phi_n(a)$ makes $a$ an element of order exactly $n$ modulo $p$, because a smaller order would make $a$ a double root of $X^n-1$, impossible when $p\nmid n$ (the derivative test). Then evaluate $\Phi_n$ at a multiple of $n$ and of all known such primes, so that the new prime factor is new.
\writeup (a) $\Phi_1=X-1$. If the claim holds for all proper divisors of $n$, then $\Phi_n=(X^n-1)/\prod_{d\mid n,d<n}\Phi_d$ is the quotient of a monic integer polynomial by a monic integer polynomial, and the division algorithm stays in $\Z[X]$ (each step divides by the leading coefficient $1$). The remainder is zero because the factorisation over $\C$ already shows exact divisibility; uniqueness of quotient and remainder identifies the integer quotient with $\Phi_n$. Evaluating at $0$: $-1=\prod_{d\mid n}\Phi_d(0)=-\prod_{d\mid n,d>1}\Phi_d(0)$, and by induction all factors except $\Phi_n(0)$ are $\pm1$, so $\Phi_n(0)=\pm1$.

(b) $a^n\equiv1\pmod p$ since $\Phi_n(a)\mid a^n-1$. Suppose $d=\operatorname{ord}_p(a)<n$. Then $a$ is a root modulo $p$ of $X^d-1=\prod_{e\mid d}\Phi_e$, hence of some $\Phi_e$ with $e\mid d$, $e<n$, $e\mid n$ ($\Z/p\Z$ is a field). So $a$ is a root of both $\Phi_e$ and $\Phi_n$, two distinct factors of $X^n-1$: it is a double root of $X^n-1$ modulo $p$, so it is also a root of the derivative $nX^{n-1}$ (differentiate $X^n-1=(X-a)^2Q$ in $(\Z/p\Z)[X]$). But $p\nmid n$ and $p\nmid a$: contradiction. So $\operatorname{ord}_p(a)=n$, and $n\mid p-1$.

(c) For $n=1$, the conclusion is Euclid's theorem. Let $n\ge2$ and suppose $p_1,\dots,p_r$ are all the primes $\equiv1\pmod n$. Let $m=np_1\cdots p_r$. For a large integer $t$, $|\Phi_n(mt)|>1$ (a non-constant polynomial tends to infinity), so it has a prime factor $q$. Since $\Phi_n(mt)\equiv\Phi_n(0)=\pm1\pmod m$, $q$ divides neither $n$ nor any $p_i$, and $q\mid\Phi_n(mt)$. By (b), $q\equiv1\pmod n$, a new prime: contradiction.
\end{solution}

\section*{Sources and further reading}

\begin{itemize}
\item K.~Ireland and M.~Rosen, \emph{A Classical Introduction to Modern Number Theory}, 2nd ed., Springer, 1990, Chapter~4. The structure of $(\Z/n\Z)^\times$ and power residues.
\item C.~F.~Gauss, \emph{Disquisitiones Arithmeticae} (1801), Section~III; English translation by A.~A.~Clarke, Springer, 1986. Where primitive roots, indices and the generalised Wilson theorem first appear.
\item M.~R.~Murty, Artin's conjecture for primitive roots, \emph{The Mathematical Intelligencer} 10 (1988), 59--67. A readable survey of the open problem quoted at the end of the chapter.
\item C.~Hooley, On Artin's conjecture, \emph{Journal f\"ur die reine und angewandte Mathematik} 225 (1967), 209--220. The proof under the generalised Riemann hypothesis.
\item W.~Diffie and M.~E.~Hellman, New directions in cryptography, \emph{IEEE Transactions on Information Theory} 22 (1976), 644--654. The paper that turned discrete logarithms into public-key cryptography.
\item P.~Moree, \emph{Artin's primitive root conjecture -- a survey}, Integers 12 (2012), A13. The 1927 conversation with Hasse, Artin's constant $0.3739\ldots$, and Heath-Brown's theorem (Quart.\ J.\ Math.\ 37 (1986), 27--38).
\end{itemize}
